\documentclass[times, twoside,onecolumn]{zHenriquesLab-StyleBioRxiv}

% Please give the surname of the lead author for the running footer
\leadauthor{Miles} 

\begin{document}

\title{Transient dynein interactions drive rapid centripetal convergence of chromosomes in early prometaphase}
\shorttitle{Early prometaphase movement}

% Use letters for affiliations, numbers to show equal authorship (if applicable) and to indicate the corresponding author
\author[1,\,\Letter]{Christopher Miles}

\author[2]{Fioranna Renda}
\author[2]{Irina Tikhonenko}
\author[2]{Alexey Khodjakov}
\author[3]{Alex Mogilner} 

\affil[1]{Department of Mathematics, University of California, Irvine, CA. }
\affil[2]{Wadsworth Center, New York State Department of Health, Albany, NY, USA.}
\affil[3]{Departments of Mathematics and Biology, New York University, NY.}


\maketitle

%TC:break Abstract
%the command above serves to have a word count for the abstract
\begin{abstract}
For proper chromosome segregation during mitosis, chromosomes must be integrated into the mitotic spindle and transported to its surface during prometaphase. By combining microscopy data with mathematical and statistical modeling, we quantify the timing and spatial dependence of early chromosome movements from their three-dimensional trajectories. We find that chromosomes start centripetal convergence to the spindle axis within just tens of seconds after the nuclear envelope breakdown. The convergence is telescopic, with the centripetal velocity as an increasing function of the distance to the spindle, resulting in chromosomes starting farther from the spindle arriving at the spindle at approximately the same time as more proximal chromosomes. Testing of mechanistic models for chromosome interactions with microtubules suggests that transient chromosome interactions with spindle microtubules are responsible for their synchronous arrival to the spindle surface. Analysis of disruption of dynein  on kinetochores suggests that this motor is responsible for the rapid chromosome capture and movements, but also that other motors account for robust but slower chromosome convergence. 
\end{abstract}
%TC:break main
%the command above serves to have a word count for the abstract

\begin{keywords}
mitosis | mathematical modeling | biophysics | prometaphase
\end{keywords}

\begin{corrauthor}
\texttt{chris.miles@uci.edu}
\end{corrauthor}

\begin{signifstatement}
The mitotic spindle is a molecular machine responsible for remarkably rapid and precise segregation of chromosomes during cell division. In early prometaphase, chromosomes converge to the spindle where they get bioriented through mechanisms still debated. By combining microscopy and modeling, we find that the farther the chromosomes are from the spindle, the earlier they start convergence and the faster they move, arriving at the spindle in synchrony with each other. Computational modeling suggests that rapid and transient, dynein motor-mediated, interactions of chromosomes with microtubules underlie the rapid and robust prometaphase.    
\end{signifstatement}


\section*{Introduction}

The mitotic spindle is a complex chemo-mechanical machine responsible for chromosomes’ segregation in mitosis \cite{prosserMitoticSpindleAssembly2017}. The remarkable spatiotemporal precision and speed of mitosis are puzzling considering the complexity and stochasticity of the spindle \cite{renda2021role}. Disentangling molecular components and mechanisms of the spindle has attracted significant combined efforts of experimental, statistical, and mathematical approaches \cite{nedelec2002computer,mogilner2010towards,redemann2019current,Miles2022} resulting in great advancements in the understanding and quantification of later phases of mitosis. 

One of the earliest phases, the onset of prometaphase, remains under-quantified and incomplete in understanding \cite{ferreira2021prometaphase}. After the NEB (nuclear envelope breakdown), for the process to continue, chromosomes (CHs) must not only be connected to the spindle itself rapidly but also transported to the spindle and bioriented (stretched and aligned with the spindle axis), so that sister chromatids are attached to the opposite spindle poles \cite{mcintoshBiophysicsMitosis2012}. A longstanding hypothesis for this rapid attachment was "search-and-capture" \cite{holy1994dynamic,heald2015thirty}, a conceptual model where microtubules (MTs) growing and shrinking from centrosomes stochastically capture CHs. 

Since its conjecture, the search-and-capture hypothesis has run into several shortcomings in reconciliation with experimental observations \cite{kliuchnikov2022celldynamo}. Even with a variety of error-correction methods \cite{zaytsev2015basic}, including biased searching \cite{wollman2005efficient}, tension sensing \cite{chen2021tension}, and geometric effects \cite{paul2009computer,Mogilner2011a,magidson2015adaptive}, the accurate integration of chromosomes into the spindle is faster and much more accurate than search-and capture predicts \cite{Renda2022,soares2022mitosis}. Notably, chromosomes get bioriented almost in sync with each other on the spindle surface, rather than at random times and places \cite{Renda2022}. The biorientation process on the spindle surface is likely driven by short MTs at the kinetochores (KTs) \cite{sikirzhytski2018microtubules}, rapid, stochastic and transient interactions of which with spindle interpolar MTs results in accurate sorting of the short MTs so that each sister KT ends up connected to the opposite polarity interpolar MTs \cite{Renda2022}. 

This leaves open the question of how most CHs are rapidly attached to the spindle and begin moving toward it arriving at the spindle surface almost simultaneously \cite{Renda2022}, considering that they are initially, at NEB, at widely varying distances from the centrosomes. One proposed mechanism is that an initial stable attachment of a CH to a long MT growing from one spindle pole and dynein-driven transport brings this CH to this pole first, and then plus-end CenpE motors on KTs are responsible for the `congression' of the CH to spindle equator where the CH finally gets bioriented \cite{skibbensDirectionalInstabilityKinetochore1993,kapoor2006chromosomes}. However, only a very small minority of CHs get mono-oriented first and pulled to one of the spindle poles \cite{magidson2011spatial}, so they must undergo congression to the equator next. Another possibility is that the process begins with stable attachments of minus-ends of two nascent K-fibers to the oppositely oriented long interpolar MTs of the spindle, and then coupled growth of the K-fibers and dynein-driven transport of their minus-ends to the opposing spindle poles completes prometaphase \cite{khodjakovMinusendCapturePreformed2003,maiatoKinetochoredrivenFormationKinetochore2004,wadsworthPluribusUnumUniversal2004}. Is one of these mechanisms, or some other, consistent with data?

In this work, to answer this question, we quantify the spatial dependence of CH movements in early prometaphase. By analyzing individual CH trajectories, we find that CHs start moving toward the spindle exceptionally soon after NEB, and paradoxically, those farther away move sooner and faster, arriving to the spindle in synch with more spindle-proximal CHs. Dissected of the CH movements demonstrate that CHs primarily directed toward the spindle center and that their speed is an increasing function of the distance to the spindle. Comparison of the experimentally measured CH trajectories in control cells and cells, in which either dynein or CENP-E motors on KTs were inhibited, with several models of motor-MT-CH interactions showed that a model of rapid transient attachments (in contrast to stable attachments proposed in previous studies) between the dynein motors on KTs and short MTs and spindle MT network fits best with the data. This supports the hypothesis that initial contact of CHs with the spindle can be made rapidly mediated by a few short MTs on KTs, and then motors bound to KTs and short MTs drive the motion toward the spindle. This motion is rapid far from the spindle, where long MTs guide CHs centripetally to the spindle center, and slow near the spindle, where long MTs pull sister chromatids in opposite parallel directions to the pole. Intriguingly, experimental results indicate that dynein on the KTs is responsible for faster attachments and movements of CHs, while dynein on the short MTs accounts for slower motion. 


\section*{Results}


\begin{figure}[h]
    \centering
    \includegraphics[width=.675\textwidth]{Figures/fig0.ai}
    \caption{\textbf{Microcroscopy images of the mitotic spindle.} 
     Chromosome arms are oriented with random orientations at NEB (A) but rapidly become splayed outward (B), as chromosomes are gathered to the spindle surface (C).First two columns are maximum-intensity projections through the entire volume of the cell. Arrows denote the positions of the spindle
poles (centrioles, labeled via centrin-GFP). Kinetochores are 
visualized with CENP-A-GFP, and microtubules are stained with a 
monoclonal antibody (DM1A). DNA is stained with Hoechst 33343.}
    \label{fig:fig0}
\end{figure}



\subsection*{Chromosomes move shortly after NEB toward the spindle} \mbox{}

At the onset of the nuclear envelope breakdown (NEB), chromosomes (CHs) are distributed uniformly throughout a cylindrical nuclear volume, a few microns tall and about 10 microns in radius, with their arms oriented randomly (\cref{fig:fig0}A). Shortly after, in early prometaphase, CHs are converging centripetally to the spindle axis (\cref{fig:fig0}B). During this convergence, kinetochores (KTs) lead the CH arms, which become increasingly parallel to their motion and orthogonal to the spindle surface. %NEB is considered to be the first time consistently centripetal rapid KT motion is detected. 
By late prometaphase (\cref{fig:fig0}C), the KTs have all been gathered to the ring-like surface (shown in \cite{Renda2022} to be the spindle surface) with their arms oriented outward. During this process, the centrosomes that organize the spindle poles start at the centers of the disk-like flat surfaces limiting the nuclear volume at the top and bottom, and then slowly drift along these surfaces in the opposite directions. We call the centrosome-centrosome axis a spindle axis and its midpoint the center of the spindle. Two dense microtubules (MT) arrays splay radially from each chromosome along the disc-like surfaces limiting the former nuclear area from the top and bottom. While the CHs that are initially far from the center move visibly centripetally, the CHs initially near the center either do not move directionally, or even squeezed out from the spindle center to end up on the ring-like surface (wrapped around the spindle axis) of the nascent spindle.  These images indicate that centripetal forces on the CHs drive their movement in early prometaphase, and we turn to quantitative details of the CH trajectories to dissect these forces. 

\begin{figure}[h]
    \centering
    \includegraphics[width=.7\textwidth]{Figures/fig1.ai}
    \caption{\textbf{Chromosome movements and attachment times.} B: Prototypical chromosomal trajectory shown in cylindrical coordinate system $r,z$. Color (light to dark) conveys time from NEB. The chromosome approaches the spindle after attachment and then becomes bioriented. Pole positions are also shown with time from NEB corresponding to darkening color. A: Cartoon describing the cylindrical coordinate system. C: The same trajectory in B shown as a function of time, with colors remaining the same. D: Aggregate statistics over all cells ($n=13$) of the cylindrical positions of chromosomes over all cells. The solid line indicates mean and dotted are $\pm$one standard deviation. E: Statistics of attachment time (after NEB) in seconds. F: Attachment time as a function of the distance from the spindle center at NEB, showing that attachment time decreases with distance. }
    \label{fig:fig1attachstats}
\end{figure}


 A typical trajectory of a CH (from now on by CH trajectory we mean the trajectory of the center-of-mass of two sister KTs) can be seen in \cref{fig:fig1attachstats}A. The trajectory is depicted in the cylindrical coordinate system of the spindle shown in \cref{fig:fig1attachstats}BC (the spindle axis is the z-direction, and the radial distance is the distance to this axis). After NEB, for a short time, the CH shifts in random undirected directions. Quickly after, the chromosome presumably attaches to the spindle and begins centripetal movement, which we denote the "attachment time" to the spindle, with precise criteria defined in Methods. This inward movement switches drastically to a more slow motion along the surface of the spindle delineating by the interpolar MT bundles \cite{Renda2022}. By measuring the distance between sister KTs and the direction of the KT-KT vector, we previously determined the moment of biorientation of the CH \cite{Renda2022} and use the same criterion here (see Methods). This trajectory demonstrates the typical behavior that biorientation takes place very rapidly after the centripetal movement of CHs switches to the lateral motion along the spindle surface \cite{Renda2022}. 

Aggregated cylindrical values as a function of time can be seen in \cref{fig:fig1attachstats}D. With approximately 200 seconds, all chromosomes have completed their inward movement. The inward movement is in the radial direction $r$ and relatively little inward movement along the spindle axis $z$ occurs. Using the criterion for  `attachment time', this switch from undirected to directed movement occurs quickly after the NEB, on the order of tens of seconds, as is seen in \cref{fig:fig1attachstats}E. Surprisingly, the attachment time is to be lowest for chromosomes that start farther away, supporting that the mechanism driving the initial movement of CHs cannot be classical search-and-capture. That is, if the rate-limiting step were the stochastic capture by MTs growing from centrosomes, those closer to the spindle surface are captured faster. Note that for CHs that were initially closer than about 4 microns to the spindle center, the attachment time statistics are probably not very meaningful, because many of such CHs remain somewhat stationary or move slightly outward. Focusing only on the the majority of CHs that are initially outside this inner volume, the observation still persists that the attachment time is a decreasing function of the initial distance to the spindle center. 

\subsection*{Chromosome centripetal velocity increases with distance from the spindle center}  \mbox{}


\begin{figure}[h]
    \centering
    \includegraphics[width=.4\textwidth]{Figures/fig2.ai}
    \caption{\textbf{Chromosome velocities.} A: Aggregated per-frame velocity magnitude for chromosomes between \SIrange{3}{5}{\um}, split into those that started less than and greater than \SI{5}{\um} away from the spindle center, showing no significant difference with a two-sample $t$-test, $T=1.8912$ and $\mathrm{df}=17958$. B: The aggregated per-frame velocity magnitudes for positions that started greater \SI{5}{\um} away from the spindle center, split into positions that are greater than or less than \SI{5}{\um} away, showing a significant deviation from a two-sample $t$-test, $T=-32.8887$ and $\mathrm{df}=16812$.  C: Cartoon depicting the velocity decomposition into the component that is direct to the center of the spindle $V_d$ and indirect $V_i$. The directionality of $V_d$ is determined by toward the spindle or away. Directionality of $V_i$. is determined by the two poles. D: Magnitude of the per-frame velocities as a function of distance. E: Magnitude of the direct component of the velocity. F: Magnitude of the indirect component of the velocity. }
    \label{fig:fig2vd}
\end{figure}

We now focus on studying the movement between the attachment event and biorientation to understand the mechanism driving the CH convergence. We first addressed the question of whether CH movements in the spindle-centric coordinate system reflect the physics of convergence. That is, what if the CHs converge to their own center-of-mass, and pull the spindle poles into some positions defined by the CHs? In \cref{sect:supp2corr} we demonstrate that the poles’ movements are not responsive to the CH movements; likely, the spindle poles are positioned by a combination of strong and balanced cortical and interpolar forces independent of weak pole-CH forces. On the other hand, the CH center-of-mass seems to follow the spindle center.

We hypothesized that distance from the spindle center is a key determinant of the CH velocity due to the geometric nature of the forces generated by the spindle. To test that the velocity is a function of the distance to the spindle only, and not of time, we split the chromosome trajectories into those that started "far" (greater than  \SI{5}{\um}) away from the spindle center and those that started "close" and compared their velocity distributions while they were at intermediate range (\SIrange{3}{5}{\um}), seen in \cref{fig:fig2vd}A. We find that in this intermediate region, the chromosomes that started far and close had no significant difference in their observed velocities. However, when looking only at those starting far away, their velocity did significantly differ in the far and close regions, seen in \cref{fig:fig2vd}B. This suggests that the velocity is indeed distance-dependent and not time-dependent. This also demonstrates that crowding of CHs close to the spindle surface at the end of the convergence process is not the cause of the slowing down of the CHs near the spindle: at the end of the convergence, CHs, which started far away from the spindle, move near the spindle in the very crowded environment with the same velocities, with which CHs initially proximal to the spindle moved at the same distances earlier on.

To understand further the distance-dependence of the velocity of chromosomes, we visualized the overall magnitude of the velocity as a function of the distance from the center of the spindle, $d$. In \cref{fig:fig2vd}C we see that the velocity magnitude increases with the distance away from the center for distances greater than 4 microns. The velocity is constant for smaller distances. To understand the directionality of this velocity, we decomposed the overall velocity into two components: the "direct" velocity along the vector toward the center of the spindle, denoted $V_d$, and the "indirect" velocity orthogonal toward the center of the spindle $V_i$. While vector projection is not typically directional, we can consider toward the spindle positive and away negative for $V_d$. For $V_i$ we denote movement primarily toward pole 1 (arbitrary) one positive and pole 2 negative. In \cref{fig:fig2vd}E we see that $V_d$ is constant below 4 microns and increases  with distance above that threshold but $V_i$ remains roughly zero on average. Thus, chromosome motion is directed toward the spindle center and  displays balanced fluctuations along the spindle axis. 

To complement this derivation of the CH velocity as a function of distance from the raw data, and to gain an initial perspective on the mechanistic nature of the distance-dependent velocity, we turn to advanced statistical theory to learn from the data what is the best-fit model for the interaction between the CHs and poles.  Namely, we examine three models: first, the observed centripetal CH convergence to the spindle center could be the result of linear superposition of movements toward each of the poles. Second, it could be that the spindle MT network is geometrically organized by the spindle in a way that integrated movements converge to the spindle center directly. The second model could be considered as a special case of the first, but for now, we ask the question of whether the superposition of movements to two poles better fits the data than movement to the spindle center. Third, it is in principle possible that CHs interact with each other, perhaps through motors mediating interactions between short MTs and KTs, and this effectively causes distance-dependent CH-CH attraction bringing all CHs to their center-of-mass, in addition to the spindle center attracting each CH. 

In \cref{sect:supp3learned} we report results from statistical learning of effective distance-dependent interactions between CHs, spindle poles, and center from the data on multiple CH trajectories. Two conclusions from this statistical learning exercise are: first, the model with CH-CH interactions predicts complex, non-monotonic, and unphysical distance dependence of the interactions. Thus, we rule out this model and conclude that the CH-CH interactions in the early prometaphase are insignificant (aside from obvious short-scale steric repulsion). Second, the model of superposition of the movement to the poles fits the data to a satisfactory degree but does not reproduce the pronounced increase of the velocity with the distance above 4 microns distance to the center. This could be an artifact from the data from the parts of CH trajectories close to the spindle center and from pole-pole increasing separation which creates an illusion of effective pole-CH repulsion. Ultimately, the model of the distance-dependent CH-spindle center attraction reproduces the velocities very similar to the ones derived directly from the data (\cref{fig:fig2vd}E). 

\subsection*{Testing mechanistic models of chromosome-microtubule interactions} \mbox{}


\begin{figure}[h]
    \centering
    \includegraphics[width=.675\textwidth]{Figures/fig3.ai}
    \caption{\textbf{Mechanistic models.} A: Cartoon depicting model 1 where attachments are stable and the distance to each pole gets reduced at a constant rate, causing chromosome movement. B: Illustrative trajectories for model 1 colored by instantaneous velocity. C: The direct and indirect components of the velocity as a function of distance to the center of the spindle. D: Cartoon depicting a model where chromosome movement is driven by transient directional forces along the microtubule lattice. E: Same as b for model 2. F: Same as c for model 2. G: Same model as d but now with polar ejection forces that repel chromosomes near poles. h, i: Same as panels above for other models. J: Another variant of model 2 where instead of forces being directed toward the poles, they are now along curved MTs. K,I: Same as panels above for other models. }
    \label{fig:fig3toymodels}
\end{figure}

We next turn our attention toward conceptual mechanistic models to explain the observed distance-dependent convergence velocity. In the first model (\cref{fig:fig3toymodels}A), long straight MTs from both poles attach stably to the sister KTs. We stress that the attachment is lateral, not between the MT plus-ends and the KTs, and it is not necessary that MTs from the opposite poles attach to different sister KTs. Also, more than a pair of MTs can connect to the CH. By stability of the attachments, we mean that the KTs do not disconnect from the MTs, but can glide along the MTs, with a constant speed, toward the MT minus-ends at the poles. Then forces driving motion effectively `reel in' these attachments at a constant rate and drive the CHs toward the spindle surface similar to a person able to lift herself to a ceiling by pulling on two ropes tied to two different points on the ceiling. This motion relies on the ability of MTs to pivot at the centrosome. In \cref{fig:fig3toymodels}B, we find that this causes velocity to \textit{decrease} with distance, also seen in \cref{fig:fig3toymodels}C. This can be understood by noting that far from the spindle, shorting the MT connections changes the angle of interaction very little and does not cause significant motion in the CH. Close to the spindle axis, forces on two rigid MTs bring the CH to the axis very rapidly. This is the prediction opposite to what we observe in experimental trajectories.

The acceleration of CHs near the spindle predicted by this model could in principle be reversed by a distance-dependent polar ejection force generated by MTs pushing on the CH arms directly by polymerization forces or through interactions with chromokinesins on the arms \cite{riederMotileKinetochoresPolar1994,tokai-nishizumiChromokinesinKidRequired2005}. Such force is supposed to be greater near the poles because a greater number of MTs would be reaching proximal CH arms. In \cref{sect:supp4sims} we explore a variant of this model with polar ejection forces but find qualitatively similar results to the model without such forces. Thus, we reject the stable attachment model. 

In \cref{fig:fig3toymodels}D we consider a different model where the CH transiently interacts with the spindle MT network. We envision that forces of constant magnitude are randomly exerted on the short timescale of 0.1 to 1 sec, the timescale of individual motor run before detachment \cite{kunwar2011mechanical}. In this model, the spindle MT network geometry is organized as if the two MT subpopulations consist of long straight centrosomal MTs radiating from each pole outward, but all that the model requires is that in the local CH vicinity there are MT lateral surfaces that guide the CH equally in the directions of two poles.  From studying this model, we find that its prediction qualitatively agrees with the experimental data and statistical model and has an increase in velocity as a function of distance (\cref{fig:fig3toymodels}e,f). This is because far away from the spindle, the accumulation of small random pulls are primarily directed toward the spindle center, while near the spindle the pulls are along its axis.   

We test how this model performs if we add the distance-dependent polar ejection forces (\cref{fig:fig3toymodels}G), and find, as expected, that this addition, while not reversing the telescopic character of the centripetal velocity (\cref{fig:fig3toymodels}I), changes the trajectories by focusing them very sharply into the spindle center (\cref{fig:fig3toymodels}H). This centering effect is not observed in the experiment (\cref{fig:fig1attachstats}D). Thus, we experimentally tested whether inhibition of chromokinesins alters the centripetal CH convergence velocity. We found that this inhibition does not change the rate of CH convergence to the spindle \cref{sect:supp1chromo}. We conclude that the polar ejection force is insignificant in early prometaphase.

Lastly, we note that the transient interaction model (\cref{fig:fig3toymodels}G), though predicting qualitatively correct distance dependence of the convergence velocity, does not accurately describe the abrupt stop of the telescopic convergence at 4 microns away from the spindle center. We hypothesize that a more accurate MT network geometry, namely MTs curving along the ellipsoidal surfaces (\cref{fig:fig3toymodels}J), as observed in spindle images, could do a better job. Indeed, in (\cref{fig:fig3toymodels}KL), we see that curved MTs cause the CHs to stop farther away from the spindle center, which aligns with the experimental observations. We tested the increasing CH-MT interaction times and found that as far as those are below a few seconds, the model predicts the centripetal telescopic velocities. Longer interactions lead to wrong predictions. In conclusion, we posit that this model of transient interactions of CHs with the spindle MT network explains the observed motion of the chromosomes.

\subsection*{Dynein inhibition makes connections and centripetal convergence slower}  \mbox{}


\begin{figure}[h]
    \centering
    \includegraphics[width=.85\textwidth]{Figures/fig4.ai}
    \caption{\textbf{Treated cells}. A: Comparison of attachment times for treated cells showing that wild-type and CENP-E have similar attachment times (no significance via two-sample $t$, $T= -0.9713$ $\mathrm{df}= 824$ but wild-type differs from both dynein inhibited ($T=-8.3593, \mathrm{df}=1119$) and double inhibited ($T=-7.1312,\mathrm{df}=822$). attachment times are later for both dynein-inhibited treated groups. 
    B: Learned statistical model for treated cells. All velocities increase with distance similar to wild-type, but CENP-E inhibited cells are slightly slower at the periphery, and dynein and double inhibited are dramatically slower. C: Conceptual model for the results depicting two types of chromosomal capture: one by long MTs attaching to dyneins at the end of kinetochore MTs, and another by kinetochore dynein. }
    \label{fig:fig4treated}
\end{figure}

If the conceptual model thus far is correct, the only obvious molecular candidate for the minus-end-directed drive is dynein. Inhibiting the ability of dynein to exert transient forces along the MT lattice should cause the motion of the CHs to cease. We find the situation to be more subtle. In \cref{fig:fig4treated}A, we see that in cells where dynein on KTs is inhibited (RPE1 Rod${}^{\Delta/\Delta}$ and GSK/Rod), the attachment times are significantly longer than control or those with CenpE motor on KTs inhibition (GSK). The convergent motion of these former two also differs from the latter two, as seen in \cref{fig:fig4treated}B. In GSK cells, the centripetal is reduced, but still significantly faster away from the spindle center. In Rod and GSK/Rod, rapid movement toward the spindles are nearly entirely gone, but slow movements remain. Altogether, this shapes the conceptual model seen in \cref{fig:fig4treated}C. We posit that the CH centripetal movements are driven by the combination of two different types of dynein interaction. In the first, non-KT dyneins at the ends of short MTs associated with KTs facilitate initial connections and subsequent movement toward the spindle, but those connections and movements are slower than the connections and motion driven by dyneins on the KT. Inhibiting CenpE likely disorganizes the short MTs near KTs which impedes only the slow dyneins on these short MTs therefore not changing the attachment times and velocities much. In rod and Rod/GSK cells, KT dynein is inhibited, and this main driver of rapid attachment and convergence is unable to act. Notably, the importance of two dynein sub-populations reverses in the process of biorientation \cite{Renda2022}: inhibition of dyneins on KTs has little effect on speed and accuracy of biorientation of the CHs that arrive to the spindle surface. 

\section*{Discussion}

To summarize, by analyzing the chromosome trajectories we found that CHs start moving toward the center of the spindle within tens of seconds after NEB, indicating rapid connection of KTs to the spindle. These attachments lead the centripetal movement toward the spindle surface, with CH arms trailing behind, to spindle MTs. This connection time does not change if CENP-E motor on KTs is inhibited and is increased two-fold if dynein motor on KTs is inhibited. The connection time shortens with increasing initial CH distance to the spindle center. CHs that are initially closer than about 4 microns to the spindle center do not converge to the center, while more distal CHs converge to the center centripetally, with velocity increasing with the distance to the center. This convergent movement slows down slightly yet remains centripetal when CENP-E is inhibited and slows down significantly when dyneins on KTs are inhibited.

We tested several mechanistic models that could explain the centripetal character of the CH convergence velocity. First, gradual slowing down of the convergence closer to the spindle center due to denser CH `crowd’ near the center is not the cause of the observed effect. Second, mutual CH-CH attraction (which are in principle possible \cite{maassInterchromosomalInteractionsGenomic2019,bradyChromosomeInteractionDistance2015} through short MT-motor interactions) cannot account for the observed velocities. Third, stable `gliding’ on intersections of long centrosomal MTs, which `reels’ the MTs in and brings the CHs to the spindle center, predicts the wrong spatial distribution of the velocities. The model that works is rapid transient CH interactions with multiple MTs of the spindle network assuming frequent, small and random displacements to both poles toward MT minus ends on the poles. Note, that because of the dynamic character of the spindle MTs, stable CH-MT interactions far from the spindle center may not be possible at all, while transient ones are easier to envision.

Interestingly, a combination of a steady centripetal motor-driven transport and opposing action of the distance-dependent polar ejection force can be ruled out, both theoretically, and experimentally, by observing the unchanged convergence speed in chromokinesin-inhibited cells. Distance-dependent polar ejection forces \cite{riederMotileKinetochoresPolar1994,tokai-nishizumiChromokinesinKidRequired2005}, but in prometaphase they could be counter-productive slowing down the CH convergence to the spindle. 

Based on the dynein on KT inhibition results, we propose the model shown in \cref{fig:fig4treated}C: dyneins drive the CHs toward the center by making rapid, short, and random, minus-end-directed excursions on spindle MTs that radiate from both poles. (Even small force of 0.1 pN, an order of magnitude less than even one dynein motor can generate, is sufficient to drive a CH across the spindle \cite{nicklasCHROMOSOMEVELOCITYMITOSIS1965,taylorBrownianSaltatoryMovements1963}.) The overall effect is like the `reeling in' long astral MTs mechanism, but transient motor-MT interactions make the convergence velocities accelerate farther from the spindle. One potentially important function of the resulting telescopic velocity is synchronized arrival of CHs onto the spindle, regardless of their initial position, which could `equalize’ the time CHs spend in metaphase as well as their chances for proper biorientation. We hypothesize that dynein on KTs is responsible for effectively faster centripetal movements and that its action is complemented by the redundant slower transport by dynein motors on the minus ends of the short MTs that connect to the KTs with their plus ends (posited in \cite{magidson2011spatial}) (\cref{fig:fig4treated}C). The result of CenpE inhibition could be explained by the role of that motor in the proper organization of the short MTs near the KTs. The double action of dyneins could add robustness to the CH convergence, while short MTs allow for greater number of contacts between the KTs and long spindle MTs. The transient character of the frequent dynein-MT interactions also ensures both utilizing all spindle MTs in local vicinity of the CH, and relatively symmetric pulling toward both poles resulting in symmetric convergence to the center regardless of the initial attachment. These results also fit into the broader recent theme of how weak, transient interactions can be surprisingly impactful drivers of motion at the molecular scale in other contexts \cite{newby2017blueprint}.

The model we propose places directionality and density requirements on the MT network of the spindle. It is now increasingly apparent that distinct populations of MTs in the spindle exist \cite{tipton2022more}, and whether this complex network of MTs \cite{kiewisz2021three} and KTs \cite{conway2021self} in the prometaphase spindle corresponds to these requirements remains to be seen. Other interesting questions are how MT branching \cite{gouveia2022acentrosomal}, viscoelastic nonlinear mechanics of MT network \cite{suresh2022modeling} and chromosome arms \cite{meijering2022nonlinear} could affect the CH convergence. Also, in other contexts, MT configurations can be optimized for cargo capture \cite{mogre2021optimizing}, and it would be interesting to investigate the optimal configurations of MTs for rapid chromosome arrival to the spindle. 

 MT pivoting has been shown to play a vital role in spindle assembly and maintenance \cite{kalinina2013pivoting,edelmaier2020mechanisms,winters2019pivoting} but also affects the CH convergence to the spindle. If the CH attachment to the long astral MTs growing from the centrosomes was stable, then dynein-driven CH gliding on these MTs toward the poles would cause pivoting of these MTs into an antiparallel configuration characteristic for the interpolar MTs delineating the spindle surface. A conceptually similar mechanism (without CHs but with minus-end motors at the MT intersections) was proposed in \cite{winters2019pivoting}. Even with transient CH-MT interactions, some effective pivoting of the MTs of the spindle network is expected. 

Several studies suggested than contraction of actomyosin networks \cite{lenartContractileNuclearActin2005,boothContractileActomyosinNetwork2019}, similar to a fishnet, can bring CHs closer to the spindle. This is especially relevant because the telescopic contraction emerges naturally in the actomyosin networks \cite{linsmeierDisorderedActomyosinNetworks2016}. Such a mechanism would also provide an attractive hypothesis why distal CHs start moving centripetally earlier than the CHs initially close to the spindle center: a peripheral contracting actomyosin network would sweep with it the distal CHs first. However, dynein inhibition results, plus the fact that KTs lead the CH centripetal motion, argue that the actomyosin contraction could be a relatively minor part of the convergence process at most.

Last, but not least, what are molecular mechanisms responsible for the shockingly rapid onset of the CH convergence? The search-and-capture-type mechanism relying on dynamically unstable MTs cannot explain the first contact between the KTs and spindle MTs sooner than several minutes after NEB \cite{magidson2011spatial,magidson2015adaptive,kliuchnikov2022celldynamo}. MT pivoting can in principle accelerate such contacts, but it would still not explain why the distal CHs make the contact sooner. We favor the following idea: assuming that $\sim$ 1000 MTs splay out from each of the centrosome and considering that the initial nuclear volume is roughly a cylinder with height 4 and radius 10 microns, then at the cylindrical surface, farthest away from the spindle center, there will still be one MT per 0.5x0.5 $\mu$m$^2$ area of this surface. Several short MTs (a few hundred nm-long each \cite{sikirzhytski2018microtubules}) rapidly pivoting around each KT would intersect with one of the long centrosomal MTs in seconds, and then motors on these intersections could make the contact. However, this hypothesis still would not explain why the distal CHs are connected first, and why the dynein on KTs, not the dyneins on the short MTs, are responsible for the fastest captures. This mystery remains for the future.

{\footnotesize 
\section*{Materials and Methods}


\subsection*{Cell culture, imaging, treatment}
The data analyzed in the work is taken from \cite{Renda2022}. Detailed methods can be found there. 
\textcolor{red}{Alexey and Fiory should put info here.}

\subsection*{Attachment time}

The chromosome positions $\vect{x}(t)$ at time $t$ are the midpoints of the two three-dimensional kinetochore positions. Velocities are defined to be $\vect{v}(t)=[\vect{x}(t+\Delta t)-\vect{x}(t)]/\Delta t$, where $\Delta t = 5\si{s}$. The pole center is $p_c(t)=(\vect{p}_1(t)+\vect{p}_2(t))/2$ where $p_1(t),p_2(t)$ are the 3D pole positions at frame $t$ in lab coordinates. The chromosome-to-pole vectors are then $\vect{d}_{1/2} = \vect{x}(t)-p_{1/2}(t)$ and chromosome-to-center is $\vect{d} = \vect{x}(t)-p_{c}(t)$. Then, $t_1 =  \arg\min_\tau \left \{\langle    \vect{v}(t) \cdot \vect{p}_1(t)/ ( \|\vect{v}(t)\|  \| \vect{p}_1(t)\|) \rangle_{t=[\tau,\tau+s]}>0.9\right\}$, where $\langle \cdot \rangle_{t=[\tau,\tau+s]}>0.9$ is the average over a time window, chosen to be $s=30\si{s}$ starting at $\tau$. In words, $t_1$ is the first time where the velocity of the chromosome move towards pole 1 for a sustained window of time. The quantities $t_2$ and $t_c$ are defined similarly for motion toward the second pole and pole center. Then, the attachment time is defined to be $t_a = \min\{t_1,t_2,t_c\}$. That is, the first time the chromosome demonstrates sustained movement toward either pole or the pole center. 

\subsection*{Biorientation time} As defined in \cite{Renda2022}, this is the the first time for each trajectory inter-kinetochore exceeds \SI{0.9}{\um} while cTilt (the relative angle between the sister KTs and the pole-pole axis) remains below $\pi/8$ radians for at least \SI{30}{s}. The median biorientation time for each cell is computed from its trajectories. The movement of chromosomes is analyzed from their attachment time (or NEB) until this median biorientation time, to account for heterogeneities in timing between cells.

\subsection*{Velocity decomposition} At each frame, the velocity $\vect{v}(t)$ is decomposed using the standard vector projection onto $\vect{d}$, the vector toward the pole-center, $\vect{v}_d(t)=\vect{d}(t)[\vect{v}(t) \cdot \vect{d}(t)]/\|\vect{d}(t)\|^2$. The orthogonal projection for the indirect component is then $\vect{v}_i(t) = \vect{v}(t)-\vect{v}_d(t)$. The magnitudes shown in the figures are given sign $V_d(t) = \sigma_d(t)\|v_d(t)\|$ where the sign is chosen by $\sigma_d(t) = 1$ if $\vect{v}(t)\cdot \vect{d}(t)>0$ and $-1$ otherwise. That is, positive if toward the spindle center. Similarly the magnitude $V_i(t)=\sigma_i(t)\|v_i(t)\|$ where $\sigma_i(t)=1$ if $\vect{v}(t)\cdot \vect{d}_1(t)>0$. That is, the sign is positive if moving more toward pole 1 than 2, and negative otherwise. 

\subsection*{Learned models} We utilize the general technique from \cite{Lu2019,Zhong2020}. The learned model shown in the main text assumes that the dynamics take the form  $\dd \vect{x} / \dd t = $ $=f(\vect{x}-\vect{p}_c) \vect{d}/\|\vect{d}\|$ $=f(\vect{x}-\vect{p}_c) (\vect{x}-\vect{p}_c)/\| \vect{x}-\vect{p}_c\| $. That is, the chromosomes move toward (or away) the pole center in a distance-dependent manner. The methods in the references outline a technique for statistically estimating the interaction $\hat{f}(\vect{d})$, which is shown in the figures. In  \cref{sect:supp3learned} we consider models where chromosomes interact with each of the poles and each other and compare the fits of each model through  cross-validation. 

\subsection*{Model simulations} All models are simulated in two spatial dimensions corresponding to the cylindrical coordinate system $r,z$. The initial positions are randomly sampled from true initial positions. The pole positions are fixed and taken to be the average of the pole positions during the movement window for the chromosomes, approximately $z=\pm\SI{3.9}{\um}$. Model 1 (stable attachments) is simulated by evolving the distance $r_1, r_2$ radially away from each chromosome. These distances $r_1, r_2$ are decreased at a constant rate (plus noise), and then the chromosome position is updated to be the intersection of the two circles with these new radii.  Model 2 is simulated by the Gillespie algorithm. The chromosome either moves toward pole 1 or pole 2 at velocities $\vect{v}_1, \vect{v}_2$ and then switches at rate $\alpha$. In the most basic model, $\alpha$  $\vect{v}_1$, and $\vect{v}_2$ are all constant. To implement polar ejection forces, $\vect{v}_{1,2}$ are modeled velocities that decrease with distance to each pole. To implement microtubule bending, the velocities are modified to reflect the ellipsoid shape of the spindle than directly toward each pole as in the basic model.  Further details can be found  in \cref{sect:supp4sims}. 

\mbox{}}


\begin{acknowledgements} This work was supported by the National Institutes of Health 
NIGMS grant GM130298 to A.K. and by the National Science Foundation grant DMS 1953430
to A.M.
\end{acknowledgements}

\section*{Bibliography}
\bibliography{chrom_move.bib}

\clearpage

%% You can use these special %TC: tags to ignore certain parts of the text.
%TC:ignore
%the command above ignores this section for word count
\onecolumn
\newpage


\clearpage 

%%%%%%%%%%%%%%%%%%%%%%%%%%%%%
% Supplementary Information %
%%%%%%%%%%%%%%%%%%%%%%%%%%%%%
\setcounter{equation}{0}
\setcounter{figure}{0}
\setcounter{table}{0}
\setcounter{page}{1}
\makeatletter
\renewcommand{\theequation}{S\arabic{equation}}
\renewcommand{\thefigure}{S\arabic{figure}}
%\renewcommand{\bibnumfmt}[1]{[S#1]}
%\renewcommand{\citenumfont}[1]{S#1}

\section*{Supplementary Materials}

% \tableofcontents


\captionsetup*{format=largeformat}
 




\section{Chromosome movement follows centrosome movement}
\label{sect:supp2corr}

All models considered in this manuscript consider CH movements relative to the spindle poles. There is a subtle assumption behind these models, namely that CH positions and movements do not affect poles' positions and movements. A mechanical way to look at this assumption is to assume that the poles are positioned by large forces between the poles and cell cortex, which are CH-independent. On the other hand, the forces between CHs and poles are weaker, not affecting the poles, yet determining CH positions and movements relative to the poles. One obvious prediction from this model is that the CH center-of-mass motion has to be in parallel, at any moment of time, to the vector of movement of the spindle center, and has to follow the spindle center, not the other way around. This allows us to disentangle the influence of pole-CH interactions and those with the cell otherwise, such as chromosome-cortex interactions \cite{kunda2009actin}. We indeed find that the chromosome movement is highly correlated with pole movement, and consequently reasonable to model as solely dependent on the pole positions themselves. This evidence can be seen in \cref{fig:suppfig2corr}. 

\begin{figure}[h]
    \centering
    \includegraphics[width=.85\textwidth]{Figures/suppfig2.ai}
    \caption{\textbf{Movement correlation}. A: Autocovariance between the movement of the center of poles and the center of chromosomes. Thin lines depict cells and pink is averaged over cells. The movement is correlated and chromosomes move with a slight lag of about 30 seconds. B: A depiction of the correlated movement in a single cell, showing that the pole center and chromosome center move in a similar manner.}
    \label{fig:suppfig2corr}
\end{figure}



In the left panel, the autocovariance between the movement of the center of mass of chromosomes with the center of the poles  can be seen. Each curve represents a different cell. The autocovariance is computed by taking the normalized dot product between the frame-to-frame movement. More explicitly, the centers of mass are $\bar{p}(t) \coloneqq \frac{1}{2}(p_1(t)+p_2(t))$ and $\bar{c}(t) \coloneqq \sum_{i=1}^{46}c_i(t)$, 
where $p_1,p_2$ are the pole positions and $c_i$ are the chromosomes, all in 3 spatial dimensions in the lab frame. The velocity of each center of mass is then$
\Delta \bar{p}(t) \coloneqq \frac{1}{\Delta t}\left( \bar{p}(t)-\bar{p}(t-\Delta t) \right)$ and  $\Delta \bar{c}(t) \coloneqq \frac{1}{\Delta t}\left( \bar{c}(t)-\bar{c}(t-\Delta t) \right)$. 
The autocovariance is then the average over all times, of the normalized dot product of each of these velocities, 
$$
\mathrm{autocov}(\mathrm{lag}) = \left \langle  \frac{\Delta \bar{p}(t) \cdot \Delta  \bar{c}(t+\mathrm{lag})}{\| \Delta \bar{p}(t)\| \|\Delta \bar{c}(t+\mathrm{lag})\|}\right \rangle_t.
$$
In the Figure, it can be seen that there is a high correlation between the directional movement between chromosomes and poles, with the average over all cells shown in the thick line (and standard deviation shaded). The slight lag (on the timescale of tens of seconds) suggests that chromosome movement lags pole movement slightly.  The right panel depicts the spatial trajectories in a single cell, illustrating that indeed the center of masses are visibly correlated. 



\section{Learned statistical pairwise interaction models}
\label{sect:supp3learned}

The method for inferring interactions in a pairwise interaction model is taken from \cite{Lu2019,Zhong2020}.  In all the considered models, we consider the pole positions $p_1(t)$ and $p_2(t)$ to be externally prescribed. That is, we do not model their movement but use their positions in the considered statistical models for chromosome movement. We assume that all chromosomes follow the same movement law, but note that a more detailed model may account for differences in the chromosomes themselves, for instance, due to size \cite{drpic2018chromosome}. The movement is then modeled by the following general form for the position of the $i$th chromosome, $x_i$:
$$
\frac{\dd x_i}{\dd t} = \mathcal{F}= \frac{1}{46}\sum_{j=1}^{46} F(x_i -x_j) \frac{x_i-x_j}{\|x_i-x_j\|} + \frac{1}{2}G(x_i - p_1)\frac{x_i-p_1}{\|x_i-p_1\|}+\frac{1}{2}G(x_i - p_2)\frac{x_i-p_2}{\|x_i-p_2\|}+ H(x_i-p_c)\frac{x_i-p_c}{\|x_i-p_c\|} + \sqrt{2D} \xi_i(t).
$$
The last term $\xi_i(t)$ is a standard white-noise term with components satisfying $\langle \xi_i(t),\xi_i(t+\tau)\rangle=\delta(t-\tau)$ and each $\xi_i$ is independent from another. Then, $F(\cdot)$ represents chromosome-chromosome interactions (perhaps via volume exclusion), $G(\cdot)$ represents interactions directly with the poles, and $H(\cdot)$ is the interaction with the center of the spindle $p_c$. We explore the following model possibilities:
\begin{enumerate}
    \item $F\equiv 0$, $G\neq 0$, $H \equiv 0$. Interaction with only the poles.
    \item $F \equiv 0 $, $G\equiv 0$, $H \neq 0$. Interaction with only the pole center. $\star$ (used in the main text). 
    \item $F \neq 0 $, $G\equiv 0$, $H \neq 0$. Interaction with the pole center and other chromosomes. 
\end{enumerate}
It should be noted that one inherent challenge in comparing these models is their differing complexities. Including more terms will inherently produce a better fit. However, to account for this, we perform cross-validation. That is, we perform 1-cross validation by repeatedly fitting the model leaving the trajectories of 1  cell out, and then compute the prediction error for the trajectories in that cell under the learned model. The error (or loss) considered is the mean squared error between the empirical velocities and those from the model. That is, the empirical velocity of each chromosome is
$$
\Delta x_i(t) = \frac{x(t+\Delta t)-x(t)}{\Delta t}.
$$
With superscript $m$ denoting each cell observed, the least squares error is then
$$
\mathcal{E} = \sum_{\substack{\mathrm{chroms }\, i\\\mathrm{times } \, t \\\mathrm{cells} \, m}} \| \Delta x_i^m(t) - \mathcal{F}(x_1, \ldots, x_{46}^m)\|^2.
$$
The premise of minimizing this error is to expand each unknown function into a known functional basis $\psi_b$ with $B$ elements.
$$
F(r) = \sum_{b=1}^{B} f_b \psi_b(r),\quad G(r) = \sum_{b=1}^{B} g_b \psi_b(r), \quad H(r) = \sum_{b=1}^{B} h_b \psi_b(r)
$$ 

In this work, we specifically take $\psi_b$ to be Chebyshev polynomials generated by $\psi_b(\cos \theta) = \cos(b \theta)$. The this relation yields  $\psi_0(r) = 1$, $\psi_1(r)=r$, $\psi_2(r)= 2r^2-1$ and so on. In this formulation, finding the basis coefficients that minimize the squared error is a linear problem and can be solved directly. The authors in \cite{Lu2019,Zhong2020}
 astutely point out efficient constructions of the linear system that allow for each cell to be considered in parallel before one final linear system solve. 
 

\begin{figure}[h]
    \centering
    \includegraphics[width=.9\textwidth]{Figures/suppfig3.ai}
    \caption{\textbf{Other learned models}. a: Mean squared error for the main statistical model (interaction with pole center) for varying number of basis elements. b: The learned interaction for the alternate model where chromosomes interact with the poles directly. c: Learned interactions for the model where chromosomes still interact with the spindle but also with each other. The model predicts nonphysical long-range chromosome-chromosome interaction.}
    \label{fig:suppfig3otherlearned}
\end{figure}

 The results of the fitting various models can be seen in \cref{fig:suppfig3otherlearned}. In the left panel, the number of basis elements $B$ is varied for the main models. The error shown is the average error from the leave-one-out validation on each cell. In the validation error calculation, the error is computed identically to the minimized error in the regression: it is assumed that the true position is known at each step and the error is computed at step. Intuitively, a bias-variance tradeoff appears as a function of $B$, where all methods perform optimally for some number of basis elements, approximately 5-10.  The figure in the main text shows $B=10$,  the basis size that yields the smallest error. In all figures, the error bars shown are the 5\% and 95\% quantiles ranges of fits from 100 bootstrapped data sets, where the same number of overall trajectories is fit, but with possible repeats. It is not shown in a figure, but a numerical comparison also reveals that for this range of basis elements, the model where chromosomes only interact with the pole center consistently fits better across all basis element ranges and therefore is presented in the main text.
 
The resulting interactions of fitting the alternate models are also shown in \cref{fig:suppfig3otherlearned}. In the model where chromosomes only interact with the poles (panel b), it can be seen that the prediction is qualitatively similar to the one in the main text: at long distances, poles attract chromosomes with relatively constant velocity. However, one difference is the appearance of strong repulsion at short distances to the poles. 

The results of fitting the model with chromosome interactions can be seen in panel c. The interaction with the pole center $H$ is qualitatively identical to the model without chromosome interactions. However, note that the fitted interaction between chromosomes and poles appears even at long ranges between them, a physical impossibility. One could conceive of modifying the loss function to penalize known physical constraints (such as monotonicity of the chromosome-chromosome interaction), but this is not within the scope of this current work. 


\section{Details on mechanistic movement models}
\label{sect:supp4sims}


To simulate all models, initial positions are randomly sampled from the true initial positions from experimental data. All simulations are performed in two spatial dimensions corresponding to the cylindrical coordinate frame $r$ and $z$. Although  poles separate during prometaphase, congression of chromosomes occurs before the largest extent of this, and therefore pole position is taken to be constant. The choice of pole position is the average of their positions over all cells during prometaphase. The output of the models are analyzed in the same manner as the experimental trajectories: the velocities are decomposed into the projection onto the vector toward the spindle center (and its orthogonal complement), called "direct" and "indirect" velocities in the text. For the analysis, all trajectories are downsampled to $200$ time samples to match the approximate sampling rate of experimental data.  

\subsection{Model 1 and variants}

The basis of model 1 is the assumption that the motion of chromosomes arises from stable attachments between the chromosomes and each of the poles. More explicitly, the model assumes that the chromosome is attached to (one or several) spindle microtubules and the motion arises from the connection between the chromosome and the poles shortening at a prescribed rate. The shortening can also be interpreted as the chromosomes being attached at the kinetochores, but a force drives them along the guiding spindle microtubules. 

To simulate the model, the initial position of each chromosome $x(0)$ is interpreted as the intersection of two circles corresponding to radial microtubules, each emanating from a pole. That is, radius 1 and 2 (corresponding to each pole) initialize at $r_1(0) = \|x(0)-p_1\|$ and $r_2(0)=\|x(0)-p_2\|$. Each radius then decreases with a velocity, and the motion of the chromosome follows the corresponding intersections of the two radii, as seen in \cref{fig:suppfi4simdetails}.


\begin{figure}[h]
    \centering
    \includegraphics[width=.825\textwidth]{Figures/suppfig5.ai}
    \caption{\textbf{Simulation details}. a: Depiction of simulation algorithm of model 1 in the main text. The movement of chromosomes is determined by sequential intersections of two circles. b: Depiction of the simulation of model 2 in the main text. Chromosomes move toward either of the poles and switch randomly. c: Depiction of the different directions chromosomes will move in model 2 and its variants. Blue to yellow color scheme indicates low to high total velocities.  }
    \label{fig:suppfi4simdetails}
\end{figure}

The radii are updated via
$$
\frac{\dd r_i}{\dd t} = v_r(r_i) + \sigma \sqrt{v_r(r_i)} \zeta_i(t),
$$
where, for the main text, $v_r(r_i) \equiv v_r^0$, a constant velocity and  $\zeta_i(t)$ correspond to white-noise terms. These differential equations are solved via Euler updates at each time step. In summary, the model simulation is performed by
\begin{pseudocode}[H]
     \caption{Model 1 simulation algorithm.}
     $x[0] \gets$ sampled from experimental initial positions\;
     $r_1[0] \gets \|x[0]-p_1\|$\;
     $r_2[0] \gets \|x[0]-p_2\|$\;

    \While{$ i \leq T/\Delta t$}{
    \tcc{update radii (no variants, constant $v_r$)}
        $r_1[i] \gets r_1[i-1] + \Delta t \cdot v_r^0+ \sigma\sqrt{ v_r^0 
    \cdot \Delta t \cdot \mathtt{randn}}$\;
        $r_2[i] \gets r_2[i-1] + \Delta t \cdot v_r^0+ \sigma\sqrt{ v_r^0 
    \cdot \Delta t \cdot \mathtt{randn}}$\; 
    \tcc{compute new intersection of circles at $(p_1,0)$ with radius $r_1[i]$ and $(p_2,0)$ with radius $r_2[i]$}
    $ x[i] \gets 
\frac{1}{2} ( p_1+p_2, 0)
+ \frac{r_2[i]^2 - r_2[i]^2}{2(p_1-p_2)^2} (p_2-p_1, 0) 
+ \frac{1}{2} \sqrt{ 2 \frac{r_1[i]^2+r_2[i]^2}{(p_1-p_2)^2} - \frac{(r_1[i]^2-r_2[i]^2)^2}{(p_1-p_2)^4}
- 1} ( 0, p_1-p_2 )$\;
    }
\end{pseudocode}

This noise form was chosen under the assumption that movement arises from discrete Poisson arrivals of forces at some rate. In the limit of fast events, the corresponding stochastic differential equation has mean and variance both equal, meaning the noise magnitude term should scale with the square root of the mean. T  This choice of scaling seemingly has no qualitative impact, and the choice made is only meant to display how fluctuations in the radii would affect chromosome movement under this model. The dimensionless parameter $\sigma$ therefore controls the noise magnitude and is chosen to qualitatively agree with experimental trajectories. 

The deterministic velocity component is written as a function of the distance to explore the possibility of polar ejection forces, $v_r(r) = v_r^0\left(1-\alpha e^{-r/r_0}\right)$. In this model, we consider the ejection forces to occur at a characteristic spatial scale $r_0$ and with $\alpha>1$ so these ejection forces dominate at these close ranges. The results of this variant can be seen in  \cref{fig:suppfig5model1variant} but no qualitative difference occurs in this variant. 

\begin{figure}[h]
    \centering
    \includegraphics[width=.85\textwidth]{Figures/suppfig4.ai}
    \caption{\textbf{Model 1 variant}. A: Variant of model 1 that also includes polar ejection forces. Behavior is qualitatively similar to model 1. B: Chromosomes are gathered to the spindle and curve inward, with increasing velocity. C: Velocity directly toward the spindle decreases as a function of distance and indirect remains relatively constant.  }
    \label{fig:suppfig5model1variant}
\end{figure}



\subsection{Model 2 and variants}


The basis of model 2 is the assumption that the chromosomes interact with the poles via transient pulling forces along spindle microtubules. The directionality of the forces is again assumed to be toward the poles and driven by dynein motors. 

To simulate this model, rates of transient forces to each pole are denoted $k_1$ and $k_2$. In theory, these could depend on the distance between the chromosome and pole to account for the gradient of microtubules emanating from each pole. However, we assume that chromosomes are able to interact with microtubules from each pole with equal probability, so $k_1 = k_2$. During the time between events, the chromosome moves at a distance-dependent velocity $v_i(x)$, where $i$ corresponds to the pole currently being interacted with. The overall summary of the model simulation can be seen in \cref{fig:suppfi4simdetails}. 


The simulation is performed by a stochastic Gillespie simulation, where the "reaction" to move to each pole occurs at rate $k_i$.The time until the next rate is drawn, and in between events, the pole moves at the velocity depending on which pole it is associated with. 

\begin{pseudocode}[H]
     \caption{Model 2 simulation algorithm.}
     $x(0) \gets$ sampled from experimental initial positions\;
    $k_\mathrm{total}\gets k_1+k_2$\;
    $p_1 \gets k_1/k_\mathrm{total}$\;
    \While{$ t \leq T$}{
    \tcc{choose next direction and time via Gillespie}
    $\tau \sim \exp(1/k_\mathrm{total})$\;
    $t \gets t+\tau$\;
    \uIf{$\mathtt{rand}<p_1$}{
    \tcc{move to pole 1}
    $x(t) \gets x(t)+v_1(x(t))\tau$
    }
    \Else{
    \tcc{move to pole 2}
    $x(t) \gets x(t)+ v_2(x(t))\tau$
    }
    }
\end{pseudocode}

Note that the output of the simulation is $x$ sampled at irregular times, but these positions are interpolated at the same rate $\Delta t$ to compare to other models.

In model 2 with no variants, this velocity is chosen to be 
$$
v_i^{\text{model 2}}(x) = v_0 \frac{p_i-x}{\|p_i-x\|}.
$$
That is, the velocity is constant and directed straight toward the corresponding pole. The model with poleward ejection forces takes the velocities to be 
$$
v_i^{\text{model 2b}}(x) = v_0 \left[ \left (1-\beta e^{-\|p_i-x\|/\ell}\right)\right] \frac{p_i-x}{\|p_i-x\|}.
$$
That is, the velocity decreases (and then becomes repellent) at a length scale $\ell$ in proximity to the corresponding pole. In the last variant of the model, the velocity is again constant in magnitude but directionality is modified to not be directed toward a pole but rather along directions that reflect the curvature of the microtubule. To motivate the functional form, one could consider the first form arising as the gradient flow of the density of microtubules emanating radially from the poles, or $v_i \sim -\nabla \Phi_i$, where $\Phi_i(x) = \|x-p_i\|_2$. Then, taking any other $p$-norm with even $p>2$ induces a more rectangular shape in the spindle, and we take $p=4$, resulting in 
$$
v_i^{\text{model 2c}}(x) = v_0 \frac{(x_1^5, x_2^5)}{(x_1^6+x_2^6)^{5/6}}.
$$
Since $k_1=k_2$, the resulting vector field for the fast-switching limit is simply $\dd x/ \dd t = \frac{k_1}{k_1+k_2}v_1(x)+\frac{k_2}{k_1+k_2}v_2(x)$. These vector fields are shown in \cref{fig:suppfi4simdetails} in the bottom panels. From this, it can be seen that the overall flow in the base model  is toward the spindle center, primarily in the radial direction and slightly centering in the $z$ direction. In the modified model, the motion associated with each pole (shown as "one pole") is now more spatially dependent and reflects the ellipsoidal shape of the spindle. The resulting net flow still attracts the chromosomes to the spindle but causes a significant slowdown closer to the poles. This is caused by an increase in regions with parallel microtubules and therefore no net force generation. This form is admittedly phenomenological and of natural interest to make more realistic in future studies. 


\begin{table}[h]
    \centering
    \begin{tabular}{l|l|l}
        parameter &  value  &  meaning  \\ \hline
       $\Delta t$ & 0.05 \si{\second}   & time step for model 1\\
       $T$ & 500 \si{\second} & max time for all models \\
      $v_r^0$ & 0.075 \si{\um\per\second} & velocity in model 1\\
       $\sigma$ & 0.1 & noise magnitude in model 1\\
       $\alpha$ & 7.5 & scale of repulsion (relative to attraction) in model 1 variant \\
       $r_0$ & 2 \si{\um} & spatial scale of repulsion in model 1 variant\\ 
       $k_1$,\, $k_2$ & 5 \si{\per\second} & rate of transient movements in model 2\\
       $v_0$ & 0.05 \si{\um\per\second} &  velocity in model 2 \\ 
        $\beta$ & 3 & scale of repulsion (relative to attraction) in model 2 variant \\
       $\ell$ & 2.5 \si{\um} & spatial scale of repulsion in model 2 variant
    \end{tabular}
    \caption{Parameter values used in simulations of the toy models.}
    \label{tab:suppparams}
\end{table}

\section{Chromokinesin inhibition}
\label{sect:supp1chromo}


One possible source of forces or attachments driving prometaphase congression are chromokinesins. These motors are known to play a role in the ejection of chromosome arms \cite{magidson2015adaptive} but longstanding evidence suggests they are not responsible for congression \cite{levesque2001chromokinesin,kline2004depletion}. To verify this, we examined the spatial patterns of chromosomes in Kid-inhibited cells and compared them with the wild-type patterns, seen in \cref{fig:suppfig1chromo}. 

\begin{figure}[h]
    \centering
    \includegraphics[width=.6\textwidth]{Figures/suppfig1.ai}
    \caption{\textbf{Chromokinesin inhibition}. Panels depict the chromosome positions in cylindrical coordinates for wild-type and chromokinesin-treated populations. Standard deviation is over a single cell. Error bars are over multiple cells.}
    \label{fig:suppfig1chromo}
\end{figure}

The figure shows a comparison of the cylindrical coordinate system $z, r$ for each cell type. Each curve denotes the average over different cells ($n=14$ for Kid-inhibited cells) with mean and standard deviation referring to within-cell spatial variation. In the figure, it can be seen that the congression of chromosomes in both control and Kid-treated cells is primarily in the radial direction, and does not seem to differ significantly. This therefore supports the suggestion that chromokinesins do not play a fundamental role in this phase of mitosis. 






%TC:endignore
%the command above ignores this section for word count

\end{document}
